% page 385 du fac-simile (image p-384 du scan)
% bloc 172.7 x 242.0 mm ; marge gauche 29.3 mm
\pgc{7.620mm}{0.000mm}{
\l{\cel{56}377}
\l{}
\l{\cel{4}mito.\cel{2}Rakonto tradicionita qua atribuas,ad ula eventi o per-}
\l{\cel{7}sonegi, karaktero supernatura. - DEFIRS.}
\l{}
\l{\cel{4}miteno. Ganto qua kovras la manuo, tote o duime, kun fingruyo}
\l{\cel{7}nur por la polexo. - EFS.}
\l{}
\l{\cel{4}\sou{mitokondrio}. (\sou{biol}.)Grani o filamenti, tre mikra, kolorizebla}
\l{\cel{7}per alizarino e la kristalo violea, qui existas ye quanti}
\l{\cel{7}plu o min granda en omna celuli, di animali o di vejetanti,}
\l{\cel{7}dum lia tota vivo. - DEFIRS.}
\l{}
\l{\cel{4}\sou{mitologio}.Historio fablatra di la dei, di la enti supernatura.}
\l{\cel{7}- DEFIRS.}
\l{}
\l{\cel{4}\sou{mitomo}. (\sou{biol.}) Parto di la citoplasmo, ek filamenti e fibreti,}
\l{\cel{7}branchigita od anostomozigita, ma ne retikuligita, qua}
\l{\cel{7}segun-dice konsistas ye protoplasmo vera.}
\l{}
\l{\cel{4}\sou{mitoso}. (\sou{biol.}) Partigeso nemediata di la celulo, en qua la}
\l{\cel{7}divido di la nukleo eventas ante olta di la korpo celula.}
\l{}
\l{\cel{4}\sou{mitro}. I. (\sou{epoki antiqua)}.\cel{2}Kapo-vesto Aziana, quan pose +weris}
\l{\cel{7}la Romanini alta-ranga. - II. Kapo-vesto di la episkopi,}
\l{\cel{7}kardinali ed ula abadi kande li oficias solene. - DEFIRS.}
\l{}
\l{\cel{4}\sou{mitralio}. I. Peci de olda fero, kupro, bronzo. - II. (\sou{olim}).}
\l{\cel{7}Olda ferajo per qua onu charjis la kanoni. - EFIS.}
\l{}
\l{\cel{4}\sou{mitralioso.}\cel{2}(\sou{milit}.)\cel{2}Kanono ye pafado rapida ed iterata. DEFR.}
\l{}
\l{\cel{4}\sou{mixar}. (\sou{trans., kun}). I. Unionar ensemble (kozi diversa) tale ke}
\l{\cel{7}la diverseso cesas esar videbla. - II. Enduktar aden medio}
\l{\cel{7}(ulu od ulo stranjera). - EFI.}
\l{}
\l{\cel{4}\sou{mixamebo}. (\sou{biol.}) Celulo qua devenas de la jermifo di la sporo}
\l{\cel{7}che la fungi mixomiceta, qua konsistas ye amaso de proto-}
\l{\cel{7}plasmo, sen membrano celuloza, e ye movi amiboidala. -}
\l{}
\l{\cel{4}mizantropo. Individuo qua odias la genero \textquotedbl{}homo\textquotedbl{}. - DEFIS.}
\l{}
\l{\cel{4}mizerar. - (\sou{netrans}.) Esar tante povra ke (lu) esas kompatinda.}
\l{\cel{7}DEFIRS.}
\l{}
\l{\cel{4}mizerikordio. Kompato qua pardonas la kulpinto. - FIS.}
\l{}
\l{\cel{4}mnemoniko. Arto helpar memoro per procedi qui igas la kozi}
\l{\cel{7}plu facile retenebla, merkebla, ritrovebla. - DEFIRS.}
\l{}
\l{\cel{4}mobilizar. (\sou{milito)}.\cel{2}Igar armeo, od armeo-fragmento, departo-}
\l{\cel{7}pronta por militesko. - DEFIRS.}
\l{}
\l{\cel{4}moblo. Objekto uzata por garnisar, ornar la interna parto di}
\l{\cel{7}domo (stuli, tabli, ed objekti diversa altra-sorta : liti,}
\l{\cel{7}e c.) - DFIRS.}
}
